% 由输入分布与核构造联合分布及输出分布。教学数据：mu=(0.5,0.3,0.2)，
% K 同 cond-finite-kernel；联合概率 M(i,j)=mu_i K(i,j)：
% (0.35,0.10,0.05)、(0.06,0.18,0.06)、(0.02,0.06,0.12)；列和 muK=(0.43,0.34,0.23)。
\begin{tikzpicture}[line width=.6pt,font=\fontsize{8}{10}\selectfont,>=Stealth,
  cell/.style={draw=FigureStroke,minimum width=11mm,minimum height=6.5mm,inner sep=0pt},
  inp/.style={cell,fill=FigureBlueFill},
  outc/.style={cell,fill=FigureGreenFill},
  note/.style={FigureMuted,font=\fontsize{7.5}{9}\selectfont}]
% 输入分布列
\node[FigureBlue] at (0,7mm) {$\mu(s)$};
\foreach \i/\k/\m in {0/1/0.5,1/2/0.3,2/3/0.2}{
  \node[FigureInk,anchor=east] at (-7mm,-6.5mm*\i) {$s=\k$};
  \node[inp] at (0,-6.5mm*\i) {$\m$};}
\node[FigureInk] at (11mm,-6.5mm) {$\times$};
\node[note,align=center,text width=12mm] at (11mm,-12mm) {乘$K$的对应行};
% 联合表
\begin{scope}[xshift=22mm]
\foreach \t/\j in {a/0,b/1,c/2}{\node[FigureInk] at (11mm*\j+5.5mm,7mm) {$t=\t$};}
\foreach \i/\a/\fa/\b/\fb/\c/\fc in {0/0.35/70/0.10/20/0.05/10,1/0.06/12/0.18/36/0.06/12,2/0.02/4/0.06/12/0.12/24}{
  \node[cell,fill=FigureBlue!\fa!white] at (5.5mm,-6.5mm*\i) {$\a$};
  \node[cell,fill=FigureBlue!\fb!white] at (16.5mm,-6.5mm*\i) {$\b$};
  \node[cell,fill=FigureBlue!\fc!white] at (27.5mm,-6.5mm*\i) {$\c$};}
\draw[->,FigureGreen] (16.5mm,-16.4mm) -- (16.5mm,-19.6mm);
\node[note,anchor=west] at (34mm,-18mm) {按列求和};
\foreach \j/\v in {0/0.43,1/0.34,2/0.23}{\node[outc] at (11mm*\j+5.5mm,-24mm) {$\v$};}
\node[FigureGreen,anchor=west] at (34mm,-24mm) {$\mu K(t)$};
\node[FigureInk] at (16.5mm,-33mm) {联合概率$M(s,t)=\mu(s)K(s,t)$};
\end{scope}
\end{tikzpicture}
